\label{sec:methods}

The framework connects three operations: defining the enclosure to be annotated, measuring each submitted layout, and combining observations into a collective layout. Controlled team comparisons then isolate changes due to team size and membership. Each submission is an ordered ring of upper--lower corner pairs. The lower endpoint determines a floor-level wall junction; the upper endpoint determines its local height. Small endpoint shifts can change the reconstructed wall substantially (Figure~\ref{fig:observed-geometry}), so both methods retain the geometry carried by these observations.

\begin{figure}[t]
\centering
\setcounter{subfigure}{0}
\begin{minipage}{82mm}
\centering
\ESPanel{.64\linewidth}{panels/fig03/a.pdf}{Panorama}%
\hfill%
\ESPanel{.33\linewidth}{panels/fig03/b.pdf}{Floor corner}

\par\vspace{3pt}
\ESPanel{.34\linewidth}{panels/fig03/c.pdf}{Full BEV}%
\hfill%
\ESPanel{.476\linewidth}{panels/fig03/d.pdf}{3D overlay}

\par\vspace{2pt}
\includegraphics[width=.64\linewidth]{\FigureAssetRoot/panels/fig03/legend.pdf}
\par\vspace{2pt}

\end{minipage}
\caption{Image observations and reconstructed geometry. A floor endpoint displacement from P to P' changes the footprint and the corresponding wall in three dimensions. The assessed layout is blue and the reference is dashed gray.}
\label{fig:observed-geometry}
\end{figure}


\subsection{Geometric quality measurement}
\label{sec:quality}

A layout can cover the correct floor area while distorting the ceiling, or have orthogonal walls around the wrong enclosure. We introduce $Q$ to measure these differences through floor overlap, boundary displacement, internal direction, mean height and local height. Figure~\ref{fig:quality-components} shows two controlled changes and their resulting scores. Geometric responses and human judgments are evaluated in Section~\ref{sec:q-validation}.

\begin{figure*}[t]
\centering
\setcounter{subfigure}{0}
\begin{minipage}{\textwidth}
\centering
\ESPanel{.24\linewidth}{panels/fig04/fig04_a.pdf}{Boundary shift: 0.4 m}%
\hfill%
\ESPanel{.24\linewidth}{panels/fig04/fig04_b.pdf}{Boundary shift: 1.2 m}%
\hfill%
\ESPanel{.24\linewidth}{panels/fig04/fig04_c.pdf}{Local height shift: $\pm 0.3$ m}%
\hfill%
\ESPanel{.24\linewidth}{panels/fig04/fig04_d.pdf}{Local height shift: $\pm 1.2$ m}

\par\vspace{2pt}
\includegraphics[width=.82\linewidth]{\FigureAssetRoot/panels/fig04/fig04_legend.pdf}
\par\vspace{2pt}

\end{minipage}
\caption{Controlled examples show how geometric changes affect Q. The reference room has a $4\times4$ m footprint, a 3 m wall height and a 1.6 m camera height. (a--b) Moving one floor boundary inward reduces overlap and increases boundary distance; hatched regions mark omitted floor area. (c--d) Opposing upper-boundary shifts preserve the footprint and perimeter-mean height while increasing local height error. Blue and orange show the smaller and larger changes; dashed gray denotes the reference. Measurements and Q are computed from the displayed geometry.}
\label{fig:quality-components}
\end{figure*}


Let $P$ and $R$ denote a layout and its assigned reference, reconstructed in camera-height units with floor plane $y=-1$. Each upper--lower pair is horizontally aligned to the midpoint of its shorter arc, retaining both vertical coordinates and hence its local wall height (Supplement B5). For floor polygons $A_P$ and $A_R$, overlap is $I=|A_P\cap A_R|/|A_P\cup A_R|$. Boundary distance $S_b$, in degrees, symmetrically averages the nearest angular distance from uniformly spaced samples on each spherical floor contour to the other continuous boundary.

Internal orthogonal regularity depends on $P$ itself. For floor edges with lengths $\ell_e$ and plan-view orientations $\alpha_e$, we fit the orthogonal axis pair that minimizes the length-weighted residual:
\begin{equation}
 D_{\mathrm{dir}}(P)=\min_{0\leq\phi<90^\circ}
 \sqrt{\frac{\sum_e\ell_e\operatorname{wrap}_{90^\circ}(\alpha_e-\phi)^2}{\sum_e\ell_e}},
 \label{eq:direction}
\end{equation}
where $\operatorname{wrap}_{90^\circ}$ maps deviations into $[-45^\circ,45^\circ)$.

Heights are interpolated linearly along floor edges. For cyclic corner heights $h_i$, the perimeter-weighted mean is
\begin{equation}
 \mu_P=\frac{\sum_i\ell_i(h_i+h_{i+1})}{2\sum_i\ell_i}.
 \label{eq:perimeter-height}
\end{equation}
Mean-height error is $H_m=|\mu_P-\mu_R|/\mu_R$. Local-height error $H_c$ is the symmetric root-mean-square difference between mean-centered height profiles, normalized by $\mu_R$. Each perimeter sample is compared with the nearest floor position on the other layout; distance ties use the largest squared height residual. Integrating along the perimeter makes both height measurements invariant to subdividing an unchanged straight wall with collinear points (exact expressions in Supplement C18).

For a valid layout with positive overlap, the measurements give
\begin{align}
 F&=\left(\frac1I-1\right)+\frac{S_b^2}{32}, &
 L_d&=1.5\frac{D_{\mathrm{dir}}^2}{D_{\mathrm{dir}}^2+8^2},\label{eq:floor-loss}\\
 L_m&=\left(\frac{H_m}{0.8}\right)^2, &
 L_c&=\left(\frac{H_c}{0.25}\right)^2,\label{eq:height-loss}
\end{align}
and
\begin{equation}
 Q(P,R)=\frac{100}{1+\sqrt{F^2+L_d^2}+L_m+L_c}.\label{eq:quality}
\end{equation}
If floor overlap is the only discrepancy, $Q=100I$. Scorable layouts have a simple, nondegenerate floor polygon, at least three corner pairs, positive wall heights and valid boundary projections. Valid polygons with zero overlap receive $Q=0$; incomplete or degenerate outputs remain unscorable and enter the completion-rate analysis (Supplement C15).

\subsection{Structural consensus from paired corner observations}
\label{sec:consensus}

A shallow recess may occupy little floor area yet introduce two wall turns that some people retain and others bridge over. Paired-corner structural consensus makes these choices explicit: it identifies corresponding pairs, counts their distinct supporters, aggregates their positions, and connects the retained pairs into a layout. Figure~\ref{fig:consensus-steps} follows actual observations into their resulting boundary and three-dimensional geometry.

\begin{figure*}[t]
\centering
\setcounter{subfigure}{0}
\begin{minipage}{170mm}
\centering
\ESPanel{.322\linewidth}{panels/fig06/a_upper_evidence.pdf}{Upper observations}%
\hfill%
\ESPanel{.322\linewidth}{panels/fig06/b_lower_evidence.pdf}{Lower observations}%
\hfill%
\ESPanel{.322\linewidth}{panels/fig06/c_actual_layout.pdf}{Connected 3D layout}
\par\vspace{2pt}
\includegraphics[width=\linewidth]{\FigureAssetRoot/panels/fig06/legend.pdf}
\par
\end{minipage}
\caption{Paired observations form the recessed wall in jtcxE69GiFV-13, Team 2. (a) Eight upper observations per turn, their medians and the resulting upper boundary. (b) The same supporters' lower observations projected to the floor, with individual wall traces in gray and consensus in purple. (c) Complete six-corner layout and enlarged upper-wall detail, using the same coordinates and view; rear edges are dashed. A and B identify the same turns throughout. Each corner retains its local height; $h$ denotes camera height.}
\label{fig:consensus-steps}
\end{figure*}


An observation $a_{wj}=(p^t_{wj},p^b_{wj})$ contains annotator $w$'s upper and lower endpoints. Matching cost is their mean spherical angular distance, $c(a,b)=[\theta(p_a^t,p_b^t)+\theta(p_a^b,p_b^b)]/2$. Dynamic programming over cyclic starting points and both traversal directions gives partial one-to-one matches while allowing unmatched corners:
\begin{equation}
 \pi_{uv}^{*}=\underset{\pi\in\Pi_{uv}}{\arg\min}
 \left[\sum_{(a,b)\in\pi}c(a,b)
 +g\bigl(n_u^{\mathrm{skip}}+n_v^{\mathrm{skip}}\bigr)\right],
 \label{eq:cyclic-match}
\end{equation}
where $n_w^{\mathrm{skip}}$ counts unmatched observations and $g$ is their per-observation skip penalty. Complete-linkage clustering partitions the pairwise matches into disjoint corner clusters $G$, requiring pairwise support throughout each cluster and at most one observation per annotator.

For team $S$, a cluster is retained when at least half the team supports it:
\begin{equation}
 z(G)=\mathrm{1}[2s(G)\geq |S|],
 \label{eq:mv50}
\end{equation}
where $s(G)$ counts distinct supporting annotators. This MV50 rule includes exact-half support. Only supporters contribute to the retained corner's coordinates, with equal weight. Horizontal positions are unwrapped around an observed anchor and aggregated by the median; upper and lower elevations use separate medians. The retained pairs are connected in cyclic wall order and checked for geometric validity (Supplements C15--C16, C18). Section~\ref{sec:method-validation} examines how support and upper endpoints enter the output.
